\clearpage
\fancyhead[L]{\small HIND'S STERN. PARTHENOPE. VICTORIA.}
\SourceLeaf{483}
\begin{center}
\rule{0.55\linewidth}{0.4pt}\\[2pt]
Astronomische Nachrichten. Band XXVIII. Nr. 650. Seite 23. 1848 November 2.\\[2pt]
\rule{0.55\linewidth}{0.4pt}
\end{center}
\begin{flushright}Göttingen, 1848 August 19.\end{flushright}

Den Meridiankreis habe ich erst Ende Juli wieder aufstellen können, und die Zeit wahrgenommen, wo es noch möglich war, den \textsc{Hind}schen Stern im Meridian zu beobachten. Meine Beobachtung gab:

\begin{center}
1848 Juli 30.\quad Ger. Aufst. $16^{\rm u}\ 51^{\rm m}\ 1^{\rm s}77$\\
Südl. Abw. $12^\circ\ 39'\ 10''6$
\end{center}

Prof. \textsc{Goldschmidt} hatte schon aus frühern den mittleren Ort bestimmt

\begin{center}
für 1848 Jan. 1.\quad Ger. Aufst. $16^{\rm u}\ 50^{\rm m}\ 59^{\rm s}06$\\
Abw. $-12^\circ\ 39'\ 17''3$
\end{center}

und hat jetzt sämmtliche uns bekannt gewordene Beobachtungen damit verglichen; \mbox{---\quad---\quad---}

\clearpage
\hypertarget{native-483-parthenope}{}
\begin{center}
\rule{0.55\linewidth}{0.4pt}\\[2pt]
Astronomische Nachrichten. Band XXX. Nr. 719. Seite 374. 1850 Juni 20.\\[2pt]
\rule{0.55\linewidth}{0.4pt}
\end{center}
\begin{center}Beobachtungen und Elemente der Parthenope.\end{center}

Herr Geh. Hofrath \textsc{Gauss} hat Parthenope am Meridiankreise der Göttinger Sternwarte beobachtet,

\begin{center}
\setlength{\tabcolsep}{4pt}
\begin{tabular}{l|r|r|r}
 & Mittlere Zeit & AR. & Südl. Abw. \\\hline
Mai 31 & $10^{\rm u}\ 27^{\rm m}\ 38^{\rm s}$ & $225^\circ\ 52'\ 38''4$ & $9^\circ\ 48'\ 3''7$ \\
\end{tabular}
\end{center}

Die Luft war in den dem Horizonte näheren Schichten etwas trübe, und unter schwächster Fadenbeleuchtung waren die Antritte nur schwer zu erkennen.

\begin{center}
\setlength{\tabcolsep}{4pt}
\begin{tabular}{l|r|r|r}
Juni 2 & $10^{\rm u}\ 18^{\rm m}\ 1^{\rm s}6$ & $225^\circ\ 30'\ 3''6$ & $9^\circ\ 45'\ 59''8$ \\
\end{tabular}
\end{center}

\clearpage
\hypertarget{native-483-victoria}{}
\begin{center}
\rule{0.55\linewidth}{0.4pt}\\[2pt]
Astronomische Nachrichten. Band XXXI. Nr. 740. Seite 305. 1850. Nov. 22.\\[2pt]
\rule{0.55\linewidth}{0.4pt}
\end{center}
\begin{flushright}Göttingen 1850 October 31.\end{flushright}

Ich übersende Ihnen, mein theuerster Freund, meine drei Meridianbeobachtungen der \emph{Victoria}, deren Schärfe, abgesehen von der Schwierigkeit des Beobachtens im, wenn auch nur wenig, beleuchteten Felde, dadurch etwas beeinträchtigt ist, dass die Reductionselemente immer erst aus entferntern Beobachtungen genommen werden konnten.

\begin{center}
\setlength{\tabcolsep}{3pt}
\begin{tabular}{l|r|r|r}
1850 & Göttinger m. Z. & Gerade Aufsteig. & Abweichung \\\hline
Sept. 30 & $10^{\rm u}\ 54^{\rm m}\ 19^{\rm s}7$ & $352^\circ\ 52'\ 43''4$ & $11^\circ\ 16'\ 48''3$ N. \\
Oct. 12 & $10\ 0\ 49.7$ & $351\ 17\ 39.6$ & $9\ 5\ 23.0$ \\
\hfill 20 & $9\ 27\ 25.9$ & $350\ 48\ 27.6$ & $7\ 45\ 38.6$ \\
\end{tabular}
\end{center}
